\documentclass[11pt]{article}
\usepackage[margin=2.5cm]{geometry}
\usepackage{amsmath,amssymb,graphicx}
\usepackage[hidelinks]{hyperref}
\title{An unmixed core in a uniform-vorticity shear layer, and a diffusion-only model of the core vorticity after roll-up}
\author{Ayumu}
\date{}
\begin{document}
\maketitle

\begin{abstract}
A shear layer that starts from a band of uniform vorticity is often cited as a case in which the Miller--Robert--Sommeria (MRS) statistical equilibrium describes the end state well. Because the initial vorticity distribution is known exactly, we compute the MRS and core--halo predictions in advance, with no adjustable parameter (for the pseudo-spectral runs, a multi-level MRS that approximates, with a finite number of levels, the distribution of the noise-free smoothed band), and compare them with the coarse-grained core vorticity of an inviscid vortex-in-cell (VIC) approximation, a hyperviscous pseudo-spectral computation that dissipates filaments, and viscous Navier--Stokes computations at four viscosities. Without viscosity the centre of the core keeps its initial vorticity and does not mix, remaining at 1.7--2.1 times the maximum vorticity predicted by MRS for the respective initial distribution, consistent with an unmixed core--halo state; this holds also when the band is halved in width so that more vortices merge. With viscosity ($\nu{}$ = 1$\times{}$10$^{-4}$ to 1$\times{}$10$^{-3}$, t $\le{}$ 400--700), the evolution of the mean core vorticity and of the energy after roll-up is reproduced by a model that keeps only diffusion (core within 0.004$\gamma{}$0 and E/E0 within 0.0013 for $\nu{}$ $\le{}$ 3$\times{}$10$^{-4}$), and the core crosses the MRS equilibrium at the same energy from above (for $\nu{}$ = 3$\times{}$10$^{-4}$ with the base edge width, at E/E0 $\approx{}$ 0.92--0.96 across three grids and up to six seeds per grid) and, within the times studied, continues below it rather than settling on it. Relaxing the viscous field while conserving its instantaneous Casimirs and energy leaves the core almost unchanged. Within the viscosities and times studied, the viscous core departs from the hyperviscous one roughly in proportion to $\nu{}$. In this set-up, agreement of the core vorticity with MRS can arise not as the destination of the fast relaxation but as a moment in the subsequent viscous decay. Over the classical Reynolds numbers and comparison times (converted from secondary sources), the difference between the core and MRS at fixed time changes sign with $\nu{}$; this does not explain the classical agreement as a whole.

\end{abstract}
\section{Introduction}

A two-dimensional inviscid flow stretches and folds its vorticity into ever finer filaments while organising into large vortices. The statistical-mechanical theory of Miller (1990) and Robert \& Sommeria (1991), hereafter MRS, predicts the end state of this process: the coarse-grained vorticity that maximises the mixing entropy subject to conservation of the energy and of all Casimirs.

Testing MRS numerically is difficult. The fine-grained vorticity distribution (the Casimirs) cannot be measured at finite resolution, so many comparisons re-extract the Casimirs from a coarse-grained field late in the integration and insert them into the theory (Dritschel et al. 2015). Long integrations usually rely on hyperviscosity; in Dritschel et al. (2015) the enstrophy decreases by 30--60\%. Modin \& Viviani (2020) asked whether this dissipation alters the long-time outcome and recomputed with a scheme that conserves the Casimirs exactly. Ryono \& Ishioka (2024) solved the full-Casimir MRS equilibrium for disordered initial fields on a sphere and found that the equilibrium vortices are much weaker and broader than those of the time integration. Mixing that stops short of equilibrium (incomplete relaxation) has been reported in simulations of decaying turbulence (Segre \& Kida 1998), and theories of isolated vortices treat it by confining the mixing to a limited region, the maximum-entropy bubble (Chavanis \& Sommeria 1998).

The shear layer, in contrast, has been regarded as a case where MRS works well. Sommeria, Staquet \& Robert (1991) integrated the Navier--Stokes equations pseudo-spectrally from a band of uniform vorticity and found the $\omega{}$--$\psi{}$ relation of the strongly mixed region in good agreement with the theory, although the flow reached a local equilibrium before mixing over the whole domain. Ryono \& Ishioka (2022) likewise showed, for a zonal initial field on a sphere, that the end state of a viscous time integration resembles an equilibrium solution. Both comparisons are made against viscous computations. A review of these applications describes the agreement of the equilibria with the observed large-scale structures in the shear layer and similar flows as qualitative, quantitative agreement having to be discussed case by case (Bouchet \& Venaille 2012). When the viscosity is small, a flow is generally thought to relax quickly to a quasi-stationary state and then to decay slowly by viscosity (Chavanis 2014). That viscosity affects the agreement with statistical-mechanical predictions has been pointed out for decaying turbulence in a closed square box. Brands et al. (1997) compared quasi-stationary states of viscous Navier--Stokes flow with an inviscid mean-field theory and related its success or failure to viscous effects, and Brands, Maassen \& Clercx (1999) found that equilibria computed from early fields did not match the quasi-stationary states, whereas those computed from late fields, close to the quasi-stationary state, did, attributing the failure to viscous dissipation and vorticity leakage at moderate Reynolds numbers. In viscous flow the vorticity--area distribution (the Casimirs) itself evolves by an anti-diffusion in vorticity space (Capel \& Pasmanter 2000). Is the agreement with MRS seen in the shear layer, a standard example of MRS working well, the destination of the first, fast relaxation, or is it produced during the second, viscous stage? For the shear layer we have not found a comparison that separates the two alongside inviscid computations.

A band of uniform vorticity is a two-level water-bag, so its fine-grained vorticity distribution is fixed exactly by the initial condition and the MRS equilibrium can be computed in advance with no adjustable parameter. In pseudo-spectral computations the edge of the band is smoothed over a few grid cells, so the initial field is not strictly two-level; the distribution of the noise-free smoothed band is nevertheless known, and an MRS equilibrium approximating it with a finite number of levels (the multi-level MRS) can be computed in advance as well. Using this property we compare, for the same initial condition, (i) an inviscid vortex-in-cell (VIC) approximation and a hyperviscous pseudo-spectral computation that dissipates filaments, (ii) pseudo-spectral Navier--Stokes computations at four viscosities and (iii) the a priori predictions of MRS and of the core--halo theory (Pakter \& Levin 2018, 2026), in terms of the coarse-grained vorticity of the core. We show three things. First, without viscosity the centre of the core keeps its initial vorticity and does not mix; it stays far from the MRS prediction, so the core after the first stage does not agree with the MRS solution (the solution of largest entropy among the branches we found). Second, with viscosity, the evolution of the core and of the energy after roll-up is reproduced by a model that switches off advection and keeps only diffusion (for $\nu{}$ $\le{}$ 3$\times{}$10$^{-4}$ the core differs by at most 0.004$\gamma{}$0 and E/E0 by at most 0.0013), and the core crosses the MRS equilibrium at the same energy (the multi-level MRS of the noise-free smoothed band of the Navier--Stokes initial field) from above (for $\nu{}$ = 3$\times{}$10$^{-4}$ with the base edge width, at E/E0 $\approx{}$ 0.92--0.96 across three grids and up to six seeds per grid; 0.916--0.928 for $\nu{}$ = 1$\times{}$10$^{-4}$--3$\times{}$10$^{-4}$ on the base grid) and moves on below it. Third, relaxing the viscous field while conserving its instantaneous Casimirs and energy leaves the mean vorticity of the core almost unchanged. As far as the core value is concerned, the agreement with MRS in the shear layer thus arises during the second, diffusive stage.

\section{Set-up and methods}

\subsection{Set-up}
We consider incompressible two-dimensional flow on the surface of a cylinder of radius R = 1. The azimuthal coordinate $\theta{}$ $\in{}$ [0, 2$\pi{}$) is periodic; in the axial direction z there are free-slip walls at z = $\pm{}Z_{\mathrm{W}}$ ($Z_{\mathrm{W}}$ = 4), where the stream function vanishes. The vorticity and stream function are related by $\nabla{}^{2}\psi{}$ = $-\omega{}$, and the velocity is u\_$\theta{}$ = $\partial{}\psi{}$/$\partial{}$z, $u_{\mathrm{z}}$ = $-\partial{}\psi{}$/$\partial{}\theta{}$.

The initial condition is a band of uniform vorticity, $\omega{}$ = $\gamma{}$0 = 1/w for |z| < w and $\omega{}$ = 0 elsewhere (total circulation 4$\pi{}$, velocity jump $\Delta{}$U = 2). The main half-width is w = 0.08 (w/L = 0.0127 for the period L = 2$\pi{}$); for comparison with Sommeria et al. (1991) we also use their thicknesses w = 0.107 (0.017L) and w = 0.214 (0.034L).

\subsection{Inviscid flow: vortex-in-cell (VIC)}
The vorticity is represented by N = 5$\times{}$10$^{5}$ particles of equal circulation, placed uniformly at random in the band; the sampling noise seeds the instability. At each step the particle circulation is distributed bilinearly onto a 256$\times{}$257 grid, $\nabla{}^{2}\psi{}$ = $-\omega{}$ is solved with a Fourier transform in $\theta{}$ and a sine transform in z, and the centred-difference velocity is interpolated bilinearly back to the particles. Time integration is fourth-order Runge--Kutta with $\Delta{}$t = 0.01. The relative change of energy up to t = 900 is 1.1$\times{}$10$^{-4}$.

\subsection{Viscous flow: pseudo-spectral Navier--Stokes}
We solve $\partial{}\omega{}$/$\partial{}$t + u$\cdot{}\nabla{}\omega{}$ = $\nu{}\nabla{}^{2}\omega{}$ (with a hyperviscous term $-\nu_{4}{}\nabla{}^{8}\omega{}$ where stated). The walls are represented by an odd image: the domain z $\in{}$ [$-Z_{\mathrm{W}}$, 3$Z_{\mathrm{W}}$) is periodic and $\omega{}$(z$\prime{}$) = $-\omega{}$(2$Z_{\mathrm{W}}$ $-$ z$\prime{}$) on z$\prime{}$ $\in{}$ [$Z_{\mathrm{W}}$, 3$Z_{\mathrm{W}}$); the odd symmetry is preserved by the dynamics and is equivalent to free-slip walls with $\psi{}$ = 0 at z = $\pm{}Z_{\mathrm{W}}$. The grid has 512 points in $\theta{}$ and 1024 in z (including the image), aliasing is removed by the 2/3 rule, and the time step is fourth-order Runge--Kutta with $\Delta{}$t = 0.005. The edge of the band is smoothed by a tanh profile of width $\varepsilon{}$ = 3dz (dz = 16/1024; $\varepsilon{}$ = 1dz for comparison), and white noise of amplitude 10$^{-3}\gamma{}$0 is added inside the band. We use $\nu{}$ = 1$\times{}$10$^{-4}$, 2$\times{}$10$^{-4}$, 3$\times{}$10$^{-4}$ and 1$\times{}$10$^{-3}$, which covers the Reynolds numbers of Sommeria et al. (1991) converted to our units (3$\times{}$10$^{-4}$ to 1.3$\times{}$10$^{-3}$).

\subsection{Equilibrium predictions}
\begin{itemize}
\item Two-level MRS. For a water-bag the coarse-grained vorticity is $\omega{}$ = $\gamma{}$0/(1 + exp($-\beta{}$($\psi{}$ $-$ $\mu{}$))). $\beta{}$ and $\mu{}$ are fixed so that the total circulation and the energy E equal their initial values (or the viscously reduced value E(t)). $\mu{}$ is found by bisection on the circulation, $\beta{}$ by bisection on log $\beta{}$, and for each $\beta{}$ the self-consistency with $\psi{}$ is solved by relaxation iteration. Branches with one core per period, a zonal state and two cores per period are computed and compared by their mixing entropy.
\item Multi-level MRS. The initial field of the pseudo-spectral computations is not two-level because of the edge smoothing (fluid of intermediate vorticity 0.05--0.95$\gamma{}$0 carries 86\% of the circulation for $\varepsilon{}$ = 3dz and 29\% for $\varepsilon{}$ = 1dz). We therefore divide the initial vorticity into K = 16 levels (10 after removing empty levels) and solve for the equilibrium that conserves the area $A_{\mathrm{k}}$ of each level $\sigma{}_{\mathrm{k}}$, $r_{\mathrm{k}}$ $\propto{}$ exp($\beta{}\sigma{}_{\mathrm{k}}\psi{}$ + $m_{\mathrm{k}}$) with $\Sigma{}_{\mathrm{k}}$ $r_{\mathrm{k}}$ = 1 at every point. The $m_{\mathrm{k}}$ are adjusted by alternating corrections of the level areas (a Sinkhorn-type iteration) and $\beta{}$ by bisection on the energy, on a 256$\times{}$513 grid. On the grid each row across the band edge carries a single vorticity value, so the number of levels saturates at the number of rows crossing the edge (K = 32 gives 11 levels; the z spacing equals that of the Navier--Stokes grid); the K = 32 solution has the same core (0.556 at r = 0.2) as K = 16, so the result does not depend on the binning; doubling the $\theta{}$ resolution to 512 (that of the Navier--Stokes grid) leaves the core, maximum and entropy unchanged to three decimals. With K = 2 and $\varepsilon{}$ $\to{}$ 0 the method recovers the two-level MRS solution to within 2\% in the core.
\item Core--halo theory (Pakter \& Levin 2018, 2026): the core keeps the vorticity $\gamma{}$0, the halo has uniform density and its outer edge is set by the separatrix.
\item None of these predictions has an adjustable parameter.
\end{itemize}

\subsection{Relaxation at the instantaneous Casimirs (relaxation equations)}
Viscosity changes the Casimirs themselves. To ask whether a viscous field at a given time is close to equilibrium for its own instantaneous Casimirs and energy, we use the relaxation equations of Robert \& Sommeria (1992) without advection. The vorticity of the Navier--Stokes field is divided into K levels ($\sigma{}_{\mathrm{k}}$, K = 10); starting from the $r_{\mathrm{k}}$ that assign each point to one level (mixed with 10$^{-3}$ of the domain-averaged level fraction to avoid the logarithmic singularity, which leaves the area and circulation of each level unchanged), we solve

  $\partial{}r_{\mathrm{k}}$/$\partial{}$t = $\nabla{}\cdot{}$[ D ( $\nabla{}r_{\mathrm{k}}$ + $\beta{}$ $r_{\mathrm{k}}$ ($\sigma{}_{\mathrm{k}}$ $-$ $\bar{\omega}$) $\gamma{}$0 $\nabla{}\psi{}$ ) ],  $\bar{\omega}$ = $\Sigma{}_{\mathrm{k}}$ $\sigma{}_{\mathrm{k}}$ $r_{\mathrm{k}}$,  $\nabla{}^{2}\psi{}$ = $-\gamma{}$0 $\bar{\omega}$.

The conservative form keeps the area and circulation of every level exactly, and $\beta{}$ is found at every step by a secant iteration so that the discrete energy change vanishes. The drift velocity |$\beta{}$($\sigma{}_{\mathrm{k}}$ $-$ $\bar{\omega}$)$\gamma{}$0$\nabla{}\psi{}$| gives a cell P\'eclet number above 10, at which upwind differencing adds more numerical diffusion than the physical D and the evolution is no longer a gradient flow. We therefore use Scharfetter--Gummel face fluxes, F = (D/$\Delta{}$)[B($-$Pe) $r_{\mathrm{i}}$ $-$ B(Pe) $r_{i+1}$] with B(x) = x/(e$^{x}$ $-$ 1) and Pe = v$\Delta{}$/D, and redistribute the mismatch of $\Sigma{}_{\mathrm{k}}$ F among the faces in proportion to the level fractions so that $\Sigma{}_{\mathrm{k}}$ $r_{\mathrm{k}}$ = 1 is kept. As checks, the one-core MRS solution (two levels) is a fixed point (core 0.455; S changes by 10$^{-5}$ in 1000 steps), and the VIC field (two levels) relaxes to almost the MRS shape (core 0.441 against 0.455; the VIC field differs by 3.6\% in energy and 1.4\% in the area of the high level, so entropy values are not compared). The relaxation increases the entropy monotonically, but where it ends depends on the starting point (reaching a local maximum is not guaranteed in general), and it is not a tool for finding the global maximum (Ryono \& Ishioka 2022 also note that the relaxation method cannot break the symmetry of the initial field).

\subsection{Reproducibility}
All computations use our own Python codes (VIC, pseudo-spectral Navier--Stokes, two-level and multi-level MRS, relaxation equations with Scharfetter--Gummel fluxes, diagnostics). The random seed is 1 for both VIC and Navier--Stokes (torch.manual\_seed). Rerunning the Navier--Stokes case $\nu{}$ = 3$\times{}$10$^{-4}$ with seed 2 changes the core by at most 0.012$\gamma{}$0 and E/E0 by at most 0.0016 for t $\le{}$ 400, and moves the crossing of the multi-level MRS(E) from E/E0 = 0.928 (t = 163) to 0.934 (t = 146). The Navier--Stokes initial field is the tanh band of section 2.3 plus white noise of amplitude 10$^{-3}\gamma{}$0. The levels of the multi-level MRS are built from the noise-free tanh band placed on the MRS grid (256$\times{}$513, with the same z spacing as the Navier--Stokes grid); the noise contributes of order 10$^{-3}$ to the circulation and its effect on the level areas is neglected in this comparison.

In the doubled-grid run (1024$\times{}$2048) the edge smoothing is set in grid points ($\varepsilon{}$ = 3dz), so its physical width is halved and E0 differs by 0.3\%. To separate the grid from $\varepsilon{}$ we also ran $\nu{}$ = 3$\times{}$10$^{-4}$ on a half grid (256$\times{}$512, $\varepsilon{}$ = 1.5dz) and a doubled grid (1024$\times{}$2048, $\varepsilon{}$ = 6dz), both with the same physical edge width as the base grid (E0 agrees to 2$\times{}$10$^{-5}$). With white noise drawn independently at every grid point, the three grids start from different initial fields, and compared at the same energy their cores differ by up to 0.05$\gamma{}$0 (at E/E0 = 0.95: 0.496, 0.519 and 0.473 on the half, base and doubled grids; crossings of the multi-level MRS(E) at E/E0 = 0.938, 0.928 and 0.951). To separate the discretisation from the initial noise we generated the noise on the half grid and interpolated it spectrally onto the finer grids (exact at the coarse grid points), so that all three grids start from the same physical field. The cores then agree to within 0.003$\gamma{}$0 between the half and base grids and to within 0.001$\gamma{}$0 between the base and doubled grids for t $\ge{}$ 40, and E/E0 to within 1$\times{}$10$^{-5}$ between the latter two; the crossing of MRS(E) is at E/E0 = 0.938, 0.940 and 0.940 on the three grids; the discretisation error of the core is therefore much smaller than the differences above, which come from the initial noise.

We measured the effect of the noise with six seeds on the half grid: at E/E0 = 0.95 the core is 0.494 $\pm{}$ 0.007 (standard deviation; range 0.486--0.503), and the crossing of MRS(E) lies at E/E0 = 0.934--0.946. Reducing the noise amplitude by a factor of 2 at a fixed seed moves the core by 0.002, and by a factor of 4 by $-$0.011, +0.008 and $-$0.003 for three seeds, with no systematic sign. With grid-point noise the base grid gives, over five seeds, 0.506 $\pm{}$ 0.010 at E/E0 = 0.95 (crossing at E/E0 = 0.923--0.937), 0.012 above the half-grid mean. Removing from the base-grid noise (seed 1) only the components the half grid cannot represent moves its core from 0.519 to 0.507, a shift equal to this offset, so the offset comes from the noise at scales the half grid cannot represent. On the doubled grid with grid-point noise, three seeds give 0.473, 0.511 and 0.462 (mean 0.482; crossings at E/E0 = 0.951, 0.931 and 0.957), below both distributions, although with the same initial field the doubled grid reproduces the base grid to 0.001; removing from the doubled-grid noise (seed 1) the components the base grid cannot represent leaves the core unchanged (0.473 in both), as expected for scales that viscosity damps well before roll-up ($\nu{}$k$^{2}$ ≳ 12 for k ≳ 200), and with three seeds the difference from the base-grid mean is not significant (Welch t = $-$1.5). In all cases the spread is set by the initial perturbation, not by the resolution. On every grid and every seed the core crosses MRS(E) once from above and then moves below it, and the diffusion-only model started at t = 50 reproduces the core and E/E0 (within 0.0035 and 0.0009 for all six half-grid seeds and all five base-grid seeds, within 0.0028 and 0.0005 on the doubled grid).

The core-circle mean is the plain average over grid points whose centres lie within radius r of the maximum of $\psi{}$. The energy is E = ½∫$\psi{}\omega{}$ d$\theta{}$ dz over the physical strip |z| < $Z_{\mathrm{W}}$ only ($\psi{}$ solved on the periodic image domain, $\psi{}$ = 0 at the walls), and E0 is its value at t = 0. Codes, run scripts, figure scripts and result files: https://github.com/tomo3141592653/shear-layer-unmixed-core (the numbers in this paper correspond to release v1.0; Python 3.13, PyTorch 2.10, NumPy 2.2, Matplotlib 3.10; runs/README.md maps each result file to its command line).

\subsection{Diagnostics}
The VIC particles and the gridded Navier--Stokes and theoretical fields are coarse-grained onto the same 64$\times{}$65 cells, normalised by the total circulation, and $\psi{}$ is recomputed.
\begin{itemize}
\item Core-side vorticity: the mean of $\omega{}$ conditioned on $\psi{}$ in 40 bins, averaged over the 5 bins of largest $\psi{}$ (in units of $\gamma{}$0).
\item Core-circle mean: the mean of $\omega{}$ within radius r = 0.05, 0.1, 0.2 of the maximum of $\psi{}$. For the Navier--Stokes runs it is taken directly on the computational grid (512$\times{}$1024, about 650 points in the r = 0.2 circle). For the diffusion-only model it is taken on the stored field (256$\times{}$256, or 512$\times{}$512 for the doubled grid), which agrees with the Navier--Stokes value to 0.001 at the starting time. On the doubled grid (1024$\times{}$2048) the Navier--Stokes and diffusion-model results differ by 0.001 in the core and 4$\times{}$10$^{-4}$ in E/E0, a difference no larger than at the base resolution, in an additional run that changes resolution and $\varepsilon{}$ together. The grid comparison at fixed $\varepsilon{}$ is given in section 2.6: for the same initial field the core agrees to 0.003$\gamma{}$0 between the half and base grids and to 0.001$\gamma{}$0 between the base and doubled grids.
\item Core circulation fraction: the circulation of cells with $\omega{}$/$\gamma{}$0 above a threshold (0.1, 0.2, 0.3) divided by the total.
\item Core ellipse: the aspect ratio $\lambda{}$ of the core (the region where $\omega{}$ exceeds half its maximum) is the square root of the ratio of the eigenvalues of its vorticity-weighted second-moment tensor, and $\omega{}_{\mathrm{c}}$ is its mean vorticity. The external strain e is the strain of the field obtained by removing the core and solving for $\psi{}$, evaluated at the core centroid, plus the strain of the core's own periodic images in $\theta{}$ (spacing 2$\pi{}$, n = $\pm{}$1\ldots{}$\pm{}$400) summed directly; removing the core from the periodic domain also removes its images, which must therefore be added back.
\item Isolated-vortex control: to test the measurement, an isolated vortex (a disc with a tanh edge, or a Gaussian) is placed in a periodic box (side 3, 256$^{2}$, absorbing layer near the edges), a uniform external strain u = (ex, $-$ey) is ramped up over a time of 40, and $\lambda{}$ and $\omega{}_{\mathrm{c}}$ are then measured with the same definitions and compared with the uniform-patch formula.
\item The rms difference of the conditional mean of $\omega{}$ on $\psi{}$ from the predictions is used only as a secondary measure, because once the core is smoothed by diffusion it no longer discriminates between the predictions (section 3.2).
\end{itemize}

\section{Results}

\subsection{Inviscid flow: the core does not mix, consistent with the core--halo picture}
The mean vorticity of the VIC core within r = 0.2 of the maximum of $\psi{}$ stays at 0.98--1.03$\gamma{}$0 from just after roll-up (t = 35) to t = 900. This agrees with the core--halo prediction of 1.00 and is far from the MRS prediction of 0.455 (figure 1). The a priori comparison with the core--halo theory is made for this core value only; the halo density and the position of the core boundary are not tested. VIC runs with seeds 3, 4 and 5 (to t = 100) give 0.94--1.03$\gamma{}$0 for t = 35--100. Doubling the particle number (10$^{6}$) and refining the grid to 384$\times{}$385 changes the core (r = 0.2) at t = 35, 48 and 100 by at most 0.021$\gamma{}$0 (0.961--0.998, within the seed range), so for t $\le{}$ 100 the core stays near its initial vorticity and far from the MRS solution independently of the particle number and grid; its persistence to t = 900 is observed at the base resolution only. About 44\% of the total circulation is in the core, independently of the threshold $\omega{}$/$\gamma{}$0 = 0.1--0.3 used to define it: the core has a sharp edge. At the band thicknesses of Sommeria et al. (1991) (w/L = 0.017 and 0.034) the core is likewise 1.0 at t = 48 and 100, far from the MRS predictions 0.550 and 0.797. (For r $\le{}$ 0.1 weak particle clustering and the bias of choosing the centre push the value slightly above $\gamma{}$0; we state results at r = 0.2.)

In the hyperviscous pseudo-spectral computation ($\nu_{4}{}$ = 10$^{-16}$) the centre of the core (r = 0.05) stays at 0.96$\gamma{}$0 for t = 70--400, the initial maximum after edge smoothing (0.957). Doubling the resolution to 1024$\times{}$2048 ($\varepsilon{}$ = 6dz, the same physical edge width) with $\nu_{4}{}$ = 3.9$\times{}$10$^{-21}$, compared with the base-grid run with $\nu_{4}{}$ = 10$^{-18}$ (same damping at the cut-off: $\nu_{4}{}k_{\mathrm{c}}^{8}$ = 0.72 for $k_{\mathrm{c}}$ = N\_$\theta{}$/3 on both grids), gives the same core: r = 0.2 within 0.02$\gamma{}$0 for t = 30--100 (0.814--0.830), the centre (r = 0.05) at 0.94$\gamma{}$0, and energy conserved to 0.08\%. The hyperviscous runs do not keep the initial Casimirs: by t = 50 they remove about half of the enstrophy (Z/Z0 = 0.49 and 0.53 on the base and doubled grids; 0.45 and 0.49 at t = 100), and the area of fluid with $\omega{}$/$\gamma{}$0 = 0.5--0.95 falls from 0.50--0.60 of the band area to 0.10--0.12 (0.23--0.25 at t = 100) as filaments are dissipated, with differences between the two grids of at most 0.04 in Z/Z0 and 0.014 in these area fractions. We therefore use the hyperviscous runs as a dissipative complement to VIC, and their comparison with MRS solved from the initial distribution is a comparison with a reference solution, not an equilibrium test with matching conserved quantities. VIC conserves the circulation of each particle, but the vorticity field reconstructed on a finite grid from a finite number of particles is not guaranteed to keep all the Casimirs. What we have checked is the robustness of the core mean to the particle number and grid, and by 'does not mix' we mean that the core keeps its initial vorticity at the core scale r = 0.05--0.2 and over the times observed. That the centre of the core keeps its initial vorticity even while hyperviscosity removes half of the enstrophy elsewhere supports the absence of mixing in the core. The mean within r = 0.2 is 0.848, 15\% below VIC, because the area of high vorticity is smaller (area with $\omega{}$ > 0.7$\gamma{}$0: 0.19 against 0.38--0.40). The size of the core depends on the initial edge smoothing (0.953 for $\varepsilon{}$ = 1dz), so we quote the r = 0.2 value as a range 0.85--1.0. Because the initial field of the hyperviscous run is not two-level (fluid of intermediate vorticity carries 86\% of the circulation), we recomputed MRS for that field's vorticity distribution (10 levels). Its core is 0.556 at r = 0.2 (two-level MRS: 0.455), and the hyperviscous core 0.848 is far from it as well. With the sharper initial edge ($\varepsilon{}$ = 1dz) the multi-level MRS core is 0.503 and the hyperviscous core 0.953, so the gap widens.

\subsection{Viscous flow: after roll-up the core value and energy are reproduced by diffusion alone}
For $\nu{}$ = 1$\times{}$10$^{-4}$ to 1$\times{}$10$^{-3}$ the core peaks at roll-up (t $\approx{}$ 30) and then decreases monotonically, crossing the MRS value at E0 earlier for larger $\nu{}$ (figure 1). We compare it with MRS solved at the instantaneous energy E(t), using the multi-level MRS built from the noise-free smoothed band of the Navier--Stokes initial field (section 2.4) (edge smoothed with $\varepsilon{}$ = 3dz), solved at E/E0 = 1.00, 0.95, 0.90 and 0.85 (core at r = 0.2: 0.556, 0.475, 0.391, 0.307). For $\nu{}$ $\le{}$ 3$\times{}$10$^{-4}$ the core lies 0.07--0.20 above MRS(E) just after roll-up, crosses MRS(E) from above at E/E0 = 0.916--0.928 (t = 609, 279 and 163 for $\nu{}$ = 1$\times{}$10$^{-4}$, 2$\times{}$10$^{-4}$ and 3$\times{}$10$^{-4}$; core 0.42--0.44), and then moves away below it (0.06--0.07 below at E/E0 $\approx{}$ 0.87). For $\nu{}$ = 1$\times{}$10$^{-3}$ the crossing happens during roll-up, and at t = 40 the core is already 0.11 below (figure 2). The core does not approach MRS(E) and stop; it passes through. The same holds for the sharper initial edge ($\varepsilon{}$ = 1dz; multi-level MRS core 0.503--0.263 for E/E0 = 1.00--0.85): $\nu{}$ = 3$\times{}$10$^{-4}$ crosses at t = 252, E/E0 = 0.904, and $\nu{}$ = 1$\times{}$10$^{-3}$ is already below at t = 40, while $\nu{}$ = 2$\times{}$10$^{-4}$ is still 0.017 above at t = 400 and $\nu{}$ = 1$\times{}$10$^{-4}$ 0.14 above. Compared with the two-level MRS (whose core is 0.10 lower) no crossing is seen and the core appears to approach from above; but the Navier--Stokes initial field is not two-level, and that impression comes from comparing with the wrong equilibrium. For $\nu{}$ $\le{}$ 3$\times{}$10$^{-4}$ the time to reach a given E scales as 1/$\nu{}$ ($\nu{}$t agrees between the three runs to within 9\% at E/E0 = 0.97 and within 1\% at E/E0 = 0.93--0.92), and the core trajectories nearly collapse when plotted against E.

Taking the Navier--Stokes field just after roll-up (t = 50) as the initial condition and evolving it by diffusion alone ($\partial{}\omega{}$/$\partial{}$t = $\nu{}\nabla{}^{2}\omega{}$, advection switched off) reproduces the core and the energy for t = 50--700 to within 0.004 and 0.0013 ($\nu{}$ = 1$\times{}$10$^{-4}$ and 3$\times{}$10$^{-4}$). Starting at t = 100 gives the same agreement (core within 0.0025), whereas starting the diffusion-only model from the band before roll-up (t = 0) produces no core and only about half the Navier--Stokes value (0.34 against 0.70 at t = 100). With the sharper initial edge ($\varepsilon{}$ = 1dz; the core itself is 0.08 higher at t = 50) the diffusion-only model again agrees to within 0.004 in the core and 0.0008 in E/E0 up to t = 400 for the three runs with $\nu{}$ $\le{}$ 3$\times{}$10$^{-4}$. Within the range studied, the post-roll-up evolution of the mean core vorticity and energy is reproduced by the model without advection, and the crossing of MRS(E) is one point on the way along which diffusion smooths the non-equilibrium structure formed at roll-up (we draw no conclusion about the shape of the core or other measures of mixing). This is a quantitative instance, in a system where MRS has been thought to work well, of the two-stage picture of Chavanis (2014) (fast relaxation followed by slow viscous decay).

We also checked continuity with the hyperviscous run (no ordinary viscosity). The difference from the hyperviscous core (same grid, same initial condition) is nearly proportional to $\nu{}$ for $\nu{}$ = 1$\times{}$10$^{-4}$ and 2$\times{}$10$^{-4}$ (the ratio of difference/$\nu{}$ between the two runs is within 5\% at t = 30 and 50 and 8\% at t = 100, and extrapolating the two points linearly to $\nu{}$ = 0 reaches the hyperviscous value to within 0.025). The same holds for $\varepsilon{}$ = 1dz: for t $\le{}$ 100 difference/$\nu{}$ agrees to 6\%, and extrapolating the two points linearly to $\nu{}$ = 0 reaches the hyperviscous value to within 0.015 (at t = 200 the proportionality is off by 13\%). Within the $\nu{}$ and time ranges studied the viscous core extrapolates smoothly to the value of the hyperviscous run; this does not show convergence to the Euler limit, in which the hyperviscosity and resolution would have to be varied together.

\subsection{Response of the core to relaxation at the instantaneous Casimirs}
Since viscosity changes the Casimirs, comparing with MRS at the initial Casimirs and the instantaneous energy cannot tell whether the flow is close to equilibrium at a given time. We therefore divided the $\nu{}$ = 3$\times{}$10$^{-4}$ field at t = 100--400 into 10 levels and raised its entropy with the relaxation equations of section 2.5, conserving the area and circulation of every level and the energy. At t = 300 the entropy increases monotonically from 0.12 to 3.49, but the core at r = 0.2 decreases only from 0.317 to 0.311 and stops (4$\times{}$10$^{4}$ steps). At t = 100, 200 and 400 the core decreases only from 0.510 to 0.500, 0.395 to 0.387 and 0.264 to 0.260. At every time the levels that rearrange as the entropy rises are those of the halo, and the core barely moves: under relaxation that conserves the discretised instantaneous Casimirs and energy, the mean vorticity of the core is nearly invariant. We have not established that a local maximum is reached or that it is stable. The pseudo-time of the relaxation is long enough to smooth the core (size about 0.2) but not to diffuse across the whole domain (length 8 in z), so the statement concerns the robustness of the mean core vorticity under finite-time relaxation, not the global maximum. With K = 5 levels the core value itself changes because it is rounded to the level means (0.299 at t = 300), but the relaxation again leaves it almost unchanged (0.299 to 0.298); with K = 20 the core goes from 0.321 to 0.317.

\subsection{Comparison at the classical conditions}
At the band thicknesses of Sommeria et al. (1991) (w/L = 0.017 and 0.034), their range of Reynolds numbers ($\nu{}$ = 3.3$\times{}$10$^{-4}$ to 1.3$\times{}$10$^{-3}$ in our units) and their comparison times (tU/L $\approx{}$ 8--16, t = 50 and 100 here), we compare the core mean (r = 0.2) with MRS (table 1). MRS is the multi-level solution for the noise-free smoothed band of the Navier--Stokes initial field ($\varepsilon{}$ = 3dz), computed at E/E0 = 1.00, 0.95, 0.90 and 0.85 and interpolated linearly to the Navier--Stokes energy at each time.


\begin{table}[t]\centering\caption{Core mean vorticity (r = 0.2, units of $\gamma{}$0) at t = 50 / 100 (VIC at t = 48 / 100). MRS(E): multi-level MRS interpolated to the Navier--Stokes energy at that time.}\label{tab1}\scriptsize\setlength{\tabcolsep}{3pt}\begin{tabular}{lccccccc}\hline
 & inviscid VIC & NS $\nu{}$=3.3$\times{}$10$^{-4}$ & MRS(E) & NS $\nu{}$=6.7$\times{}$10$^{-4}$ & MRS(E) & NS $\nu{}$=1.3$\times{}$10$^{-3}$ & MRS(E) \\\hline
w/L=0.017 & 1.016 / 1.029 & 0.635 / 0.564 & 0.588 / 0.554 & 0.522 / 0.450 & 0.546 / 0.484 & 0.363 / 0.297 & 0.492 / 0.402 \\
w/L=0.034 & 0.979 / 1.005 & 0.884 / 0.852 & 0.786 / 0.757 & 0.779 / 0.719 & 0.753 / 0.695 & 0.610 / 0.531 & 0.697 / 0.596 \\
\hline\end{tabular}\end{table}

For w/L = 0.017 the core lies 0.047 above MRS(E) at t = 50 and 0.010 above at t = 100 for $\nu{}$ = 3.3$\times{}$10$^{-4}$, just before crossing; for $\nu{}$ = 6.7$\times{}$10$^{-4}$ it is already 0.024 below at t = 50, and for $\nu{}$ = 1.3$\times{}$10$^{-3}$ it is 0.10--0.13 below. For w/L = 0.034 the core lies 0.10 above for $\nu{}$ = 3.3$\times{}$10$^{-4}$, 0.02 above for 6.7$\times{}$10$^{-4}$ and 0.07--0.09 below for 1.3$\times{}$10$^{-3}$, so the switch from above to below lies between 6.7$\times{}$10$^{-4}$ and 1.3$\times{}$10$^{-3}$ (the thicker band keeps a stronger core, shifting the switch to larger $\nu{}$). At both thicknesses, within the classical comparison times and range of Reynolds numbers, the difference between the core and MRS(E) at fixed time changes sign with $\nu{}$. Against the two-level MRS both $\nu{}$ = 3.3$\times{}$10$^{-4}$ and 6.7$\times{}$10$^{-4}$ appear 0.05--0.13 above, but this is the same mismatch as in section 3.2. The inviscid VIC core is 1.0 at both thicknesses, far from MRS for the same two-level initial distribution (0.550 and 0.797 at E0, section 3.1); the multi-level MRS of the Navier--Stokes initial field (0.637 and 0.827) refers to a different initial distribution and is only a secondary comparison, but is also far. This is consistent with the classical report that the theory agrees in the strongly mixed region while the whole domain does not mix; whether the core value agrees with MRS, however, depends on the $\nu{}$ and time chosen for the comparison. (The definition of the Reynolds number and U are taken from secondary sources.)

\section{Discussion}

\textbf{What the agreement with MRS consists of.} The shear layer has been a standard example of MRS working well, but in our comparison, for the core, the situations in which it appears to work arise during the second, viscous stage of diffusive decay. Without viscosity the centre of the core keeps its initial vorticity (in both VIC and the hyperviscous runs) and stays at 1.7--2.1 times the maximum vorticity predicted by MRS for the respective initial distribution (VIC: 2.1 times the two-level maximum 0.47; hyperviscous: 1.7 times the multi-level maximum 0.58). With viscosity the evolution of the core and energy after roll-up is reproduced by diffusion alone, and the core crosses MRS at the same energy once and moves away below it. Within the range of $\nu{}$ studied the difference between the viscous and hyperviscous cores is roughly proportional to $\nu{}$, with no jump between the two pictures, so that "agreement or disagreement can be produced by choosing $\nu{}$ and the time of comparison" holds continuously. For w = 0.08 and $\nu{}$ = 1$\times{}$10$^{-3}$ the crossing occurs during roll-up and the core is already 0.11 below at t = 40. MRS(E) solved at the initial Casimirs does not act as an attractor for the viscous flow; within the range studied the change of the core value is reproduced by diffusion alone (since viscosity changes the Casimirs, the comparison with equilibrium at the instantaneous Casimirs is made separately in section 3.3). This is consistent with the report of Sommeria et al. (1991) that the theory agrees in the strongly mixed region while the whole domain does not mix. Their comparison, however, concerned the $\omega{}$--$\psi{}$ relation of the strongly mixed region, not the core vorticity examined here. What we can state is that in this set-up agreement of the core with MRS can be produced by viscosity, and that this agreement is a point on the decay trajectory that diffusion alone reproduces after roll-up; we have not shown that the contribution of advection is small in general or identified the mechanism uniquely, and we have not established the mechanism of the classical agreement as a whole (we have not checked the conversion of their Reynolds number or their comparison region against the original).

\textbf{Relation to the two-stage picture.} That a flow relaxes quickly to a quasi-stationary state and then decays viscously (Chavanis 2014) is known in general terms. Our results give a concrete, quantitative instance in a system where MRS was thought to work well, showing (i) that the core value after the first stage does not agree with the MRS solution of largest entropy among the branches we found and is consistent with the core--halo picture of an unmixed core, and (ii) that the evolution of the mean core vorticity and energy in the second stage is reproduced by a diffusion-only model without advection. That nonlinear transfer nearly stops once the vorticity is single-valued along streamlines, after which a coherent vortex decays linearly under the dissipation, was observed qualitatively in hyperviscous simulations on the torus and sphere by Qi \& Marston (2014); here we test it quantitatively for viscous runs, with the MRS prediction known in advance. Mixing that stops short of equilibrium has also been reported for decaying turbulence from multi-level initial vorticity (Segre \& Kida 1998). Capel \& Pasmanter (2007) constructed viscous evolution models that conserve the "mass" of each microscopic vorticity value and showed that, if mixing is much faster than the changes in energy and masses, the quasi-stationary state takes the MRS form, explaining why MRS works better than expected in viscous experiments. In our shear layer the fast mixing ends at roll-up before the core reaches MRS (the inviscid core does not mix), and the subsequent evolution is diffusive; the fast-mixing condition is not met for the core.

\textbf{Equilibrium theory alone does not select the end state.} Ryono \& Ishioka (2022) showed, for a barotropically unstable zonal flow on a sphere, that the wavenumber-2 equilibrium resembling the end state of the time integration is a saddle of the mixing entropy, the local maximum being the wavenumber-1 solution, and attributed the selection to the linear instability of the initial field. Ryono \& Ishioka (2024) found that equilibrium vortices for disordered initial fields are weaker than those of the time integration. Our finding that the core stays stronger than MRS and unmixed points the same way: the end state is decided less by a comparison of equilibria than by the dynamics of roll-up, which determine how far mixing proceeds. From this viewpoint the relaxation equations (Robert \& Sommeria 1992) are best used not to search for the global equilibrium but to measure how the core responds to a finite-time increase of entropy under the given conservation constraints; Ryono \& Ishioka (2022) also note that the relaxation method cannot break the symmetry of the initial field.

\textbf{Why the diffusion-only model works, and its scope.} After roll-up the viscous core is, for t $\ge{}$ 100, a nearly steady ellipse in the external strain (major axis along $\theta{}$, at 45$^\circ$ to the extension axis), with $\omega{}$ nearly a function of $\psi{}$. The quantities compared (the mean within a circle around the maximum of $\psi{}$, and the energy) are hardly affected by small changes of the elliptical shape, so the diffusion-only model without advection reproduces them. The agreement of these diagnostics shows that their changes can be reproduced without advection; it does not imply that advection or mixing is absent. The shape, however, is not reproduced. The aspect ratio of the region above half maximum becomes rounder under diffusion alone ($\nu{}$ = 3$\times{}$10$^{-4}$: 1.18 to 1.10 for t = 100 to 400) but more elongated in the Navier--Stokes run (1.18 to 1.38), whose core also drifts in $\theta{}$. Indeed, advection does stretch the core. This shape change can be read as a steady elliptical vortex in an external strain (Moore \& Saffman 1971; we have not read the original and derived the balance from Kirchhoff rotation against strain, which agrees with equations (10)--(11) of Guha et al. 2013). The external strain e is the sum of two parts. The strain of the vorticity outside the core (the halo) is nearly constant for t = 50--400 (0.15--0.18); the strain of the core's own periodic images in $\theta{}$ (spacing 2$\pi{}$), summed directly, is 0.05--0.06, within 0.2\% of the point-vortex-row estimate $\Gamma{}_{\mathrm{core}}$ $\pi{}$/(6(2$\pi{}$)$^{2}$) used by Guha, Rahmani \& Lawrence (2013); and the extension axes of both lie at $-$45$^\circ$. We compare with the pure-strain case and do not include the background rotation due to the halo vorticity surrounding the core (the general case of Kida 1981). The major axis of the core lies at 0$^\circ$, satisfying the condition for a steady ellipse (45$^\circ$ from the extension axis). As the core weakens by diffusion, e/$\omega{}_{\mathrm{c}}$ including the images rises from 0.040 to 0.085 ($\nu{}$ = 3$\times{}$10$^{-4}$, t = 50 to 400) and the aspect ratio from 1.17 to 1.38. Inserted into the balance for a uniform elliptical patch, the measured e/$\omega{}_{\mathrm{c}}$ is 1.09--1.17 times the prediction ($\nu{}$ = 3$\times{}$10$^{-4}$) and 1.07--1.21 times ($\nu{}$ = 1$\times{}$10$^{-4}$), so the formula gives an aspect ratio 0.02--0.06 larger than measured; the points for both $\nu{}$ follow the same path just above the balance curve in the $\lambda{}$--e/$\omega{}_{\mathrm{c}}$ plane (figure 3). This 10--20\% difference is not due to edge smoothness or to measuring $\lambda{}$ by second moments (a smooth, thin-skirted isolated vortex placed in a uniform strain gives $\lambda{}$ within 0.003 of the formula, and a thick-skirted Gaussian vortex is if anything more elongated than the formula, 1.17 against 1.145), nor to non-uniformity of the strain over the core (the strain averaged over the core differs from the central value by 3\%). The ratio depends on the threshold used to define the core: at 0.7 of the maximum it is 1.04--1.11 for $\nu{}$ = 3$\times{}$10$^{-4}$ (and scatters over 0.89--1.24 for $\nu{}$ = 1$\times{}$10$^{-4}$), and at 0.3 it is 1.22--1.28 for both, the inner core being closer to the formula on average. The difference therefore lies within the ambiguity of where a smooth core ends. The shape of the weakening core is semi-quantitatively consistent, within the pure-strain approximation, with a quasi-static passage through a sequence of steady ellipses in the external strain (as a comparison of time scales, the Kirchhoff rotation period 2$\pi{}$($\lambda{}$+1)$^{2}$/($\omega{}_{\mathrm{c}}$ $\lambda{}$) is 4--10, much shorter than the time for the core vorticity to halve, about 300 for $\nu{}$ = 3$\times{}$10$^{-4}$ and 650 for $\nu{}$ = 1$\times{}$10$^{-4}$). Our claims are restricted to the core value and the energy.

\textbf{Relation to the roll-up of a uniform band.} Pozrikidis \& Higdon (1985) showed with contour dynamics that a band of uniform vorticity rolls up into nearly elliptical uniform patches, and Guha et al. (2013) confirmed with DNS of a smoothed band (Re = 10$^{4}$) that the elliptical core rotates and nutates according to Kida's equations. That our inviscid (VIC) core is a rotating ellipse is the same picture. What we add is that this core does not become the mixed core predicted by MRS, the route by which the viscous core settles onto a steady ellipse in the external strain, and the differences due to the numerical method and the edge smoothing.

\textbf{Limitations.} (1) The size of the core (the mean within r = 0.2) depends on the method and the initial edge smoothing, ranging over 0.85--1.0. That the centre of the core does not mix, and the discrimination from MRS, do not depend on this range. The elliptical dynamics of the core also differ between methods: in VIC the core keeps rotating in the external strain (aspect ratio oscillating between 1.11 and 1.38), whereas in the hyperviscous run it settles on a steady ellipse at 45$^\circ$ to the extension axis. Inserting the measured external strain (including the images) and core vorticity (e/$\omega{}_{\mathrm{c}}$ $\approx{}$ 0.022 for VIC) into the equations for a uniform ellipse (the pure-strain case of Kida 1981, in the form of equations (10)--(11) of Guha et al. 2013) gives a half-rotation time of 1.18 and an aspect-ratio oscillation of 1.14--1.36 (largest when the major axis lies along $\theta{}$), in agreement with the VIC rotation (measured 1.2 and 1.11--1.38); the hyperviscous aspect ratio 1.10 nearly coincides with the fixed point of the same equations (1.11--1.12 for e/$\omega{}_{\mathrm{c}}$ = 0.027--0.029). Whether an orbit rotates or librates about the fixed point is set by the pair (e/$\omega{}$, $\lambda{}$). The hyperviscous core is elongated in its interior at t = 100 (aspect ratio 1.95) and on a rotating orbit, then rounds by t = 200 (1.10) into the librating region ($\pm{}$1$^\circ$). Since e/$\omega{}$ is nearly constant meanwhile, the hyperviscous core settles at the fixed point because it becomes round, whereas the VIC core stays near 1.3 without rounding. The process by which an elongated vortex sheds filaments and becomes round may differ between methods (it may proceed more easily under hyperviscosity, which removes filaments). In the viscous runs the librating region widens as the core weakens and e/$\omega{}$ grows, and the rotating ellipse is captured into libration by t $\approx{}$ 100 ($\varepsilon{}$ = 1dz, $\nu{}$ = 1$\times{}$10$^{-4}$: a rotating orbit at t = 50 and libration of $\pm{}$10$^\circ$ at t = 100; measured orientations +38$^\circ$ and $-$11$^\circ$). In the hyperviscous run with the sharper edge ($\varepsilon{}$ = 1dz) the core reaches 0.953 at r = 0.2, close to the VIC value 1.0, its orientation keeps changing (+12$^\circ$ to $-$52$^\circ$ between snapshots), and its aspect ratio is largest when the orientation is near 0$^\circ$, as in VIC; three of the four fields at t = 100--400 lie on rotating orbits of the same equations. This outcome was predicted before the computation and supports a contribution of the initial edge smoothing to the difference between the VIC and hyperviscous methods (smoother-edged vortices axisymmetrise and settle on the fixed point). (2) The definition of the Reynolds number and the details of the classical comparison are taken from secondary sources. (3) The domain is one period of the cylinder (circumference 2$\pi{}$, about 80 band half-widths). Roll-up produces about three vortices (VIC at t = 6), which merge into a single core by t = 18, so merging happens once, just after roll-up. With half the band width (w = 0.04, walls at $\pm{}$2), which doubles the circumference relative to the band and the number of vortices at roll-up, VIC also merges into a single core by t = 18, and the core mean (r = 0.1, corresponding to r = 0.2 for w = 0.08) stays at 0.98--1.03 for t = 35--400 without mixing. The two-level MRS core for the same geometry is 0.29 (maximum 0.29), so the gap from the mixing prediction is larger than for w = 0.08. The stored times are too sparse to count the merger stages. Longer domains in which slow mergers cascade over many stages (selection of the longest wavelength, Thess et al. 1994) are not treated. (4) The MRS solutions are computed in the canonical form ($\beta{}$ fixed by bisection on E); we have not verified that they are local maxima in the microcanonical sense (they have the largest entropy among the branches found). (5) The core value of the viscous runs at a given energy depends on the initial noise. For the same initial field the discretisation error is below 0.003$\gamma{}$0 (section 2.6), but the seed gives a standard deviation of 0.007--0.010$\gamma{}$0 at E/E0 = 0.95 (six seeds on the half grid, five on the base grid), and white noise drawn at every grid point puts more small-scale perturbation on finer grids, which raises the base-grid mean by 0.012$\gamma{}$0. This is a dependence of the late state on the initial perturbation of the kind reported by Segre \& Kida (1998), whose fields with the same initial energy, enstrophy and vorticity distribution decayed to different final states. The finest scales present only on the doubled grid have no effect, and its lower three-seed mean is within the seed scatter. The crossing energy of MRS(E) therefore carries an uncertainty of about 0.035 in E/E0 (0.92--0.96 for $\nu{}$ = 3$\times{}$10$^{-4}$ with the base edge width). The statements that the core crosses MRS(E) and moves below it, and that the diffusion-only model reproduces the core after roll-up, hold on all three grids and for every seed.

\section*{Code and data}
All simulations and analyses were run with our own Python codes (vortex-in-cell, pseudo-spectral Navier--Stokes, two-level and multi-level MRS solvers, relaxation equations with Scharfetter--Gummel fluxes, diagnostics). The codes, the exact command lines of every run (\texttt{runs/}), the scripts that produce every figure, the table and the numbers quoted in the text (\texttt{paper\_numbers.py}), and the result files they read are available at https://github.com/tomo3141592653/shear-layer-unmixed-core (code MIT, result files CC BY 4.0).

\section*{References}
\begin{small}\begin{list}{}{\leftmargin=1.5em\itemindent=-1.5em\itemsep=1pt}
\item Bouchet, F. \& Venaille, A. 2012 Statistical mechanics of two-dimensional and geophysical flows. Phys. Rep. 515, 227--295.
\item Brands, H., Stulemeyer, J., Pasmanter, R. A. \& Schep, T. J. 1997 A mean field prediction of the asymptotic state of decaying 2D turbulence. Phys. Fluids 9, 2815--2817.
\item Brands, H., Maassen, S. R. \& Clercx, H. J. H. 1999 Statistical-mechanical predictions and Navier--Stokes dynamics of two-dimensional flows on a bounded domain. Phys. Rev. E 60, 2864--2874.
\item Capel, H. W. \& Pasmanter, R. A. 2007 Mixing and coherent structures in 2D viscous flows. arXiv:physics/0702040.
\item Capel, H. W. \& Pasmanter, R. A. 2000 Evolution of the vorticity-area density during the formation of coherent structures in two-dimensional flows. Phys. Fluids 12, 2514--2521.
\item Chavanis, P.-H. \& Sommeria, J. 1998 Classification of robust isolated vortices in two-dimensional hydrodynamics. J. Fluid Mech. 356, 259--296.
\item Chavanis, P.-H. 2014 Statistical mechanics of two-dimensional point vortices: relaxation equations and strong mixing limit. Eur. Phys. J. B 87, 81.
\item Dritschel, D. G., Qi, W. \& Marston, J. B. 2015 On the late-time behaviour of a bounded, inviscid two-dimensional flow. J. Fluid Mech. 783, 1--22.
\item Guha, A., Rahmani, M. \& Lawrence, G. A. 2013 Evolution of a barotropic shear layer into elliptical vortices. Phys. Rev. E 87, 013020.
\item Kida, S. 1981 Motion of an elliptic vortex in a uniform shear flow. J. Phys. Soc. Jpn. 50, 3517--3520.
\item Miller, J. 1990 Statistical mechanics of Euler equations in two dimensions. Phys. Rev. Lett. 65, 2137--2140.
\item Modin, K. \& Viviani, M. 2020 A Casimir preserving scheme for long-time simulation of spherical ideal hydrodynamics. J. Fluid Mech. 884, A22.
\item Moore, D. W. \& Saffman, P. G. 1971 Structure of a line vortex in an imposed strain. In Aircraft Wake Turbulence and Its Detection (ed. J. H. Olsen, A. Goldburg \& M. Rogers), pp. 339--354. Plenum.
\item Pakter, R. \& Levin, Y. 2018 Nonequilibrium statistical mechanics of two-dimensional vortices. Phys. Rev. Lett. 121, 020602.
\item Pakter, R. \& Levin, Y. 2026 Predicting the final stationary state of sheared 2D Euler fluids. J. Stat. Mech. 2026, 103203.
\item Pozrikidis, C. \& Higdon, J. J. L. 1985 Nonlinear Kelvin--Helmholtz instability of a finite vortex layer. J. Fluid Mech. 157, 225--263.
\item Qi, W. \& Marston, J. B. 2014 Hyperviscosity and statistical equilibria of Euler turbulence on the torus and the sphere. J. Stat. Mech. 2014, P07020.
\item Robert, R. \& Sommeria, J. 1991 Statistical equilibrium states for two-dimensional flows. J. Fluid Mech. 229, 291--310.
\item Robert, R. \& Sommeria, J. 1992 Relaxation towards a statistical equilibrium state in two-dimensional perfect fluid dynamics. Phys. Rev. Lett. 69, 2776--2779.
\item Ryono, K. \& Ishioka, K. 2022 New numerical methods for calculating statistical equilibria of two-dimensional turbulent flows, strictly based on the Miller--Robert--Sommeria theory. Fluid Dyn. Res. 54, 055505.
\item Ryono, K. \& Ishioka, K. 2024 Statistical equilibria of two-dimensional turbulent flows for generic initial vorticity fields on a sphere, calculated on the basis of the original Miller--Robert--Sommeria theory. Fluid Dyn. Res. 56, 065509.
\item Segre, E. \& Kida, S. 1998 Late states of incompressible 2D decaying vorticity fields. Fluid Dyn. Res. 23, 89--112.
\item Sommeria, J., Staquet, C. \& Robert, R. 1991 Final equilibrium state of a two-dimensional shear layer. J. Fluid Mech. 233, 661--689.
\item Thess, A., Sommeria, J. \& J\"uttner, B. 1994 Inertial organization of a two-dimensional turbulent vortex street. Phys. Fluids 6, 2417--2429.

\end{list}\end{small}
\begin{figure}[p]\centering\includegraphics[width=0.85\textwidth]{fig_core_circle.png}\caption{Mean vorticity within r = 0.2 of the core centre (maximum of $\psi{}$), in units of $\gamma{}$0, against time, for w = 0.08. Dark blue with markers: inviscid VIC. Purple: hyperviscous pseudo-spectral run ($\nu_{4}{}$ = 10$^{-16}$) from the same smoothed initial field as the Navier--Stokes runs. Green to red: Navier--Stokes with $\nu{}$ = 1$\times{}$10$^{-4}$, 2$\times{}$10$^{-4}$, 3$\times{}$10$^{-4}$ and 1$\times{}$10$^{-3}$. Horizontal lines: core--halo prediction (1.00), multi-level MRS approximating the noise-free smoothed band (0.56) and two-level MRS (0.45), all at the initial energy.}\label{fig1}\end{figure}
\begin{figure}[p]\centering\includegraphics[width=0.85\textwidth]{fig_core_energy.png}\caption{Mean core vorticity (r = 0.2) of the Navier--Stokes runs (t $\ge{}$ 40) against the energy E/E0, compared with MRS solved at the same energy. Black squares, dashed: multi-level MRS approximating the noise-free smoothed band (E/E0 = 1.00, 0.95, 0.90, 0.85). Grey squares, dotted: two-level (water-bag) MRS. Yellow dashed and dotted: $\nu{}$ = 3$\times{}$10$^{-4}$ on the half and doubled grids with the same physical edge width (section 2.6). Yellow band: range of six random seeds on the half grid (t $\le{}$ 200), the spread due to the initial noise. The cores cross the multi-level MRS from above near E/E0 $\approx{}$ 0.93 and continue below it (0.92--0.96 across grids and seeds for $\nu{}$ = 3$\times{}$10$^{-4}$, section 2.6).}\label{fig2}\end{figure}
\begin{figure}[p]\centering\includegraphics[width=0.85\textwidth]{fig_core_strain.png}\caption{Ratio of external strain to core vorticity e/$\omega{}_{\mathrm{c}}$ against the aspect ratio $\lambda{}$ of the core ($\omega{}$ above half its maximum), for $\nu{}$ = 1$\times{}$10$^{-4}$ (t = 100--700) and 3$\times{}$10$^{-4}$ (t = 100--400). The strain includes the core's own periodic images. Solid curve: steady uniform elliptical patch, e/$\omega{}$ = $\lambda{}$($\lambda{}$ $-$ 1)/(($\lambda{}$ + 1)($\lambda{}^{2}$ + 1)).}\label{fig3}\end{figure}
\begin{figure}[p]\centering\includegraphics[width=0.85\textwidth]{fig_nu_limit.png}\caption{Core mean vorticity (r = 0.2) against $\nu{}$ at t = 30, 50 and 100; $\nu{}$ = 0 is the hyperviscous run (no ordinary viscosity, $\nu_{4}{}$ = 10$^{-16}$). Lines: straight lines through the $\nu{}$ = 1$\times{}$10$^{-4}$ and 2$\times{}$10$^{-4}$ runs, extended to $\nu{}$ = 0, where they pass within 0.025 of the hyperviscous value; the departure from the hyperviscous core is roughly linear in $\nu{}$ at small $\nu{}$ ($\nu{}$ = 3$\times{}$10$^{-4}$ already lies above the line, $\nu{}$ = 1$\times{}$10$^{-3}$ far off).}\label{fig4}\end{figure}
\end{document}
